Mesin Turing ditemtokaké déning himpunan prentahé. Himpunan
iki minangka himpunan winates saka tupel lima. Kanggo saben
pasangan kahanan lan simbol kang diwaca, prentah nemtokaké
kahanan anyar, simbol kang ditulis, lan obahé sirah.
Himpunan winates saka tupel lima bisa dietung, mula saben
himpunan prentah mesin Turing bisa dipasangaké karo sawijining
wilangan. \emph{Indèks} mesin Turing yaiku wilangan kang
dipasangaké karo himpunan prentahé miturut skéma kang wis
ditetepaké. Kanthi cara iki, kita bisa ngrembug mesin Turing
lumantar indèksé.

Salah siji masalah wigati ngenani tindaké mesin Turing yaiku
apa mesin iku pungkasané mandheg. Tetepna fungsi $h(e,n)=1$
yèn mesin Turing mawa indèks~$e$ mandheg nalika diwiwiti
nganggo masukan~$n$, lan $h(e,n)=0$ yèn ora. Iki diarani
\emph{fungsi mandheg}. Pitakonan apa fungsi iku dhéwé bisa
dikomputasi déning mesin Turing diarani \emph{masalah mandheg}.
Wangsulané ora. Masalah mandheg ora bisa dirampungi. Asil iki
dibuktèkaké nganggo argumèn diagonal.

Masalah mandheg minangka salah siji tuladha saka kelas masalah
kang luwih amba, yaiku pitakonan apa $X$ bisa ditindakaké
nganggo mesin Turing. Masalah wigati liyané yaiku
\emph{masalah kaputusan kanggo logika tataran kapisan}.
Apa ana mesin Turing kang bisa mutusaké apa sembarang ukara
sah? Wangsulané uga ora. Iki dibuktèkaké nganggo argumèn
reduksi. Kanggo saben mesin Turing~$M$ lan masukan~$w$,
kita bisa gawé ukara tataran kapisan $T(M,w)\lif E(M,w)$
kang sah yèn lan mung yèn $M$ mandheg nalika diwiwiti nganggo
masukan~$w$. Yèn masalah kaputusan bisa dirampungi, cara iku
uga bisa dienggo ngrampungaké masalah mandheg.
